
 Here, in \cref{sync_max:sec:sync_max_aut}, we introduce sync-maximal automata as automata
 whose sets of synchronizing words have maximal state complexity.
 Then, in \cref{sync_max:sec:structure}, we give structural results.
 The main result is \cref{sync_maximal:thm:structure}, 
 which states that these automata contain
 a completely reachable and completely distinguishable subautomaton
 that differs from the whole automaton by at most one ``dangling'' state.
 Lastly, \cref{sync_max:sec:examples} provides 
 various examples illustrating 
 relations between sync-maximal, completely distinguishable, and completely reachable automata.
 For instance, we present sync-maximal automata that are not completely reachable, as well as completely reachable automata that are not sync-maximal, along with other examples illustrating our results.
 
 \section{Definition and First Properties}
\label{sync_max:sec:sync_max_aut}
	
%!TEX root = ./../../springer-mono-format/book.tex
%!TEX root = ./../../thesis-trier.tex

\IN{sync-maximal automaton}
Let $\semiautomat{A}{Q}{\Sigma}{\delta}$ be an $n$-state automaton.
Recall that by \THM~\ref{sync_max:thm:upper_bound_SynA} 
a minimal automaton for $\Syn(\mathcal A)$ has at most $2^n - n$ states. %(see \SEC~\ref{sync_max:sec:power_automat_for_SynA}). 
We call~$\mathcal A$ %is said to be 
\emph{sync-maximal} if the minimal automaton of $\Syn(\mathcal A)$ 
has $2^n - n$ states.

A sync-maximal automaton $\mathcal A$ with non-empty state set must be synchronizing.
For assume $\Syn(\mathcal A)$ is empty.
Then it is acceptable by a one-state automaton with empty final state set.
Now $2^n - n$ equals one only for $n \in \{0,1\}$.
So, $\mathcal A$ has a single state,
but this implies $\Syn(\mathcal A) = \Sigma^*$, contradicting the assumption.
Next, in \cref{sync_max:thm:char_sync_max_aut} 
we prove that in a sync-maximal automaton also all subsets of size at least two must be reachable.

\begin{theorem}
	\label{sync_max:thm:char_sync_max_aut}
	An automaton $\semiautomat{A}{Q}{\Sigma}{\delta}$ with non-empty~$Q$
	is sync-maximal if and only if the following two conditions
	hold:
	\begin{enumerate} 
		\item \label{sync_max:thm:char_sync_max_aut:1} $\mathcal A$ is synchronizing (i.e., at least one singleton subset
		is reachable) %in $\mathcal P(\mathcal A)$,
		and every non-empty subset with at least two states is reachable.
		% in $\mathcal P(\mathcal A)$.
		
		\item $\mathcal A$ is completely distinguishable.
	\end{enumerate}
\end{theorem}
\begin{proof}
	We can assume $|Q| \ge 2$. Let $\mathcal F$ be a maximal family of pairwise distinguishable and reachable subsets. 
	The family contains at most one singleton subset, as two singleton subsets are indistinguishable,
	and only non-empty subsets, as the empty set is not reachable. So, we have $\mathcal F \subseteq \{ A \subseteq Q : |A| \ge 2 \} \cup \{\{q\}\}$
	for some $q \in Q$.	As $| \{ A \subseteq Q : |A| \ge 2 \} \cup \{\{q\}\} | = 2^n  - n$,
	by Theorem~\ref{sync_max:thm:lower_bound_SynA} $\mathcal A$ is sync-maximal
	if and only if $\{ A \subseteq Q : |A| \ge 2 \} \cup \{\{q\}\} = \mathcal F$,
	which is equivalent to complete distinguishability, synchronizability,
	and reachability of all subsets of size at least two.
\end{proof}

Next, we give two corollaries to \cref{sync_max:thm:char_sync_max_aut}.

\begin{corollary}	\label{sync_max:cor:char_sync_max_aut}
	Let $\semiautomat{A}{Q}{\Sigma}{\delta}$.
	Then $\mathcal A$ is sync-maximal
	if and only if one of the following is true:
	\begin{enumerate}
		\item \label{cha_sync_max:item:Q1} $|Q| = 1$;
		\item \label{cha_sync_max:item:Q2} $|Q| = 2$ and~$\mathcal A$ is synchronizing;
		\item \label{cha_sync_max:item:Q3} $|Q| \ge 3$,  all subsets $A \subseteq Q$
		with $|A| \ge 2$ are reachable, and $\mathcal A$ is completely distinguishable.
	\end{enumerate}
\end{corollary}
\begin{proof}
	\ref{cha_sync_max:item:Q1}: The set of synchronizing words of a one-state automaton is $\Sigma^*$
	and a minimal automaton accepting all words has a single state.
	
	\ref{cha_sync_max:item:Q2}: By \cref{syncmax:rem:two-state-compl-dist}, every two-state automaton is completely distinguishable
	and the whole state set is reachable by the empty word.
	So, by \cref{sync_max:thm:char_sync_max_aut},
	sync-maximality reduces to the reachability of a single state, i.e., synchronizability.
	Thus, $\mathcal A$ is sync-maximal if and only if it is synchronizable.\footnote{For an alternative
		argument, observe that for a two-state automaton $\mathcal A$
		we have $\Syn(\mathcal A) = \Sigma^* \Gamma \Sigma^*$
		for $\Gamma = \{a\in \Sigma \mid \delta(Q, a) \ne Q \}$.
		So, the minimal automaton of $\Syn(\mathcal A)$ 
		has two states if and only if $\Gamma \ne \emptyset$,
		which is equivalent to $\Syn(\mathcal A) \ne \emptyset$ in this case.}
	
	\ref{cha_sync_max:item:Q3}:
	If $\mathcal A$
	is completely distinguishable, then $\mathcal A$
	is synchronizing  by \cref{sync_max:prop:dist_sync} as $|Q| \ge 3$.
	So, when also all subsets of size at least two
	are reachable, then \cref{sync_max:thm:char_sync_max_aut} 
	yields sync-maximality.\footnote{By \crefitem{sync_max:cor:k-sets-implies-sync}{sync_max:cor:2-sets-implies-sync}, reachability of all $2$-sets also implies
		synchronizability for $|Q| \ge 3$.}
	Conversely, if $\mathcal A$
	is sync-maximal, then, by \cref{sync_max:thm:char_sync_max_aut},
	all subsets of size at least two are reachable
	and $\mathcal A$ is completely distinguishable.
\end{proof}

\begin{corollary}
	\label{sync_max:cor:SCC}
	An automaton is strongly connected and sync-maximal if and only if
	it is completely reachable and 
	completely distinguishable.
\end{corollary}
\begin{proof}
	By \cref{sync_max:thm:char_sync_max_aut}
	a sync-maximal automaton is completely distinguishable.
	Furthermore, every subset of size at least two is reachable
	and some singleton subset is reachable.	
	Since in a strongly connected automaton, some singleton subset is reachable if and only if all singleton subsets are reachable,	we obtain complete reachability.
	Conversely, by \cref{prelim:lem:compl_reach_implies_SCC_sync} and 
	\cref{sync_max:thm:char_sync_max_aut}, a completely reachable
	and completely distinguishable automaton is strongly connected and sync-maximal.
\end{proof}

\Cref{sync_max:cor:SCC} together with \cref{sync_max:thm:cerny_compl_reach,sync_max:thm:cerny_compl_dist}
yield that the \Cerny automata are sync-maximal (this also follows from \cref{sync_max:thm:circular_sync_max} or from \cref{cor:defectatmostone-weaker} given later).

\begin{theorem}
	\label{sync_max:thm:cerny_sync_maximal}
	The \Cerny automata $\mathcal C_n$ from \cref{intro:fig:cerny}
	are sync-maximal.
\end{theorem}

\section{The Structure of Sync-Maximal Automata}
\label{sync_max:sec:structure}

%!TEX root = ../../thesis-trier.tex

Here, in \cref{sync_maximal:thm:structure}, 
we show that an automaton is sync-maximal if and only if it contains a completely reachable
and completely distinguishable subautomaton
that misses at most one ``dangling'' state.
Before stating this result, we give a lemma needed in its proof.

\begin{lemma}
	\label{sync_maximal:lem:n-1-upperset}
	Let $\semiautomat{A}{Q}{\Sigma}{\delta}$.
	Suppose $S \subseteq Q$
	is an upper set, $|S| = |Q| - 1$,
	and $u \in \Sigma^*$
	is a word of shortest length such that $\delta(Q, u) = S$.
	Then $u \in \Sigma$, i.e., $|u| = 1$.
\end{lemma}
\begin{proof}
	We can write $u = av$ for some $a \in \Sigma$ and $v \in \Sigma^*$
	where we must have $\delta(Q, a) \ne Q$
	by the choice of~$u$ as a shortest word.
	Since $|\delta(Q, u)| = |Q| - 1$
	we also have $|\delta(Q, a)| = |Q| - 1$.
	Write $\delta(Q, a) = Q \setminus \{p\}$ for some $p \in Q$.
	Assume $p \in S$.
	Then $\delta(p, v) \in S$.
	With $\delta(Q \setminus \{p\}, v) = S$
	this gives $\delta(Q, v) = S$, contradicting the minimal
	choice of~$u$. Thus, we have $Q = S \cup \{p\}$
	and so $\delta(Q, a) = S$.
\end{proof}

\begin{remark}
	\Cref{sync_maximal:lem:n-1-upperset} does not generalize
	in the sense that upper sets of size $|Q| - k$
	for some $1 \le k < |Q|$ are reachable
	by a word of length~$k$.
	For example, let $\semiautomat{A}{Q}{\{a,b,c\}}{\delta}$
	with $Q = \{0,1,2,3\}$,
	$a$ and~$b$ being simple idempotents
	mapping~$0$ to~$1$ and~$2$ to~$3$,
	and~$c$ being a permutational letter
	that swaps~$1$ and~$2$ and fixes
	every other state.
	Then a shortest word mapping~$Q$
	to the upper set $\{3\}$
    is, for example, $abcb$.
\end{remark}

By \cref{sync_max:cor:SCC}, complete reachability and complete distinguishability imply sync-maximality. By \cref{sync_max:thm:char_sync_max_aut}, every sync-maximal automaton is completely distinguishable. However, 
there exists sync-maximal automata that are not completely reachable, see \cref{sync_max:sec:examples}.

The next result also characterizes sync-maximal automata that are not completely reachable. It turns out that in every such automaton, exactly one singleton subset is unreachable.

\begin{theorem}
\label{sync_maximal:thm:structure}
	Let $\semiautomat{A}{Q}{\Sigma}{\delta}$
	be an automaton with at least three states
	and let
	\[
	S = \{ q \in Q \mid \exists u \in \Sigma^* : \applytrans{\delta}{Q}{u} = \{q\}\}
	\]
	be the set of synchronizing states (which forms a sink component by \cref{prelim:thm:SCC_sync_words}).
	Then $\mathcal A$ is sync-maximal if and only if 
	one of the following is true:
	\begin{enumerate}
		\item \label{sync_maximal:thm:structure:1} $\mathcal A$ is completely reachable and completely distinguishable or
		
		\item \label{sync_maximal:thm:structure:item:2nd} $Q = S \cup \{q\}$ for some $q \notin S$
		and, with $\Gamma = \{ x \in \Sigma \mid \applytrans{\delta}{q}{x} = q\}$, we have	
		\begin{enumerate}
			\item \label{sync_maximal:thm:structure:item:defone}
			$\applytrans{\delta}{Q}{a} = S$ for some $a \in \Sigma \setminus \Gamma$,
					
            \item \label{sync_maximal:thm:structure:item:AB} \label{sync_maximal:thm:structure:item:dupl-set} 
            for each $p \in Q$, all pairs of non-empty, non-singleton subsets
            of $\delta^{-1}(\dupl{a}, a) \cup \{p\}$ that are both distinct from
            $\delta^{-1}(\dupl{a}, a)$ and contain~$q$ are distinguishable,\footnote{Let $\{s,t\} = \delta^{-1}(\dupl{a}, a)$. This condition is equivalent
            	to distinguishability of $\{s,t,p\}$
            	and $\{p,q\}$ if $q \in \{s,t\}$;
            	and to pairwise 
            	distinguishability
            	of $\{s,t, q\}$, $\{s, q\}$,
            	and $\{t, q\}$ 
            	if $q \notin \delta^{-1}(\dupl{a}, a)$.
            	Note that for $q \notin \{s,t\}$
            	we can assume $p = q$ since otherwise no
            	subset of $\{s,t,p\}$ contains~$q$.}
                     				
			\item \label{sync_maximal:thm:structure:item:complreach} $\restrictAutGamS{\mathcal A}{\Gamma}{S}$ is completely reachable,
			
			\item \label{sync_maximal:thm:structure:item:compldist} $\mathcal A_{|S}$ is completely distinguishable.			
		\end{enumerate}
	\end{enumerate}
\end{theorem} 
\begin{proof}	
	First, suppose $\mathcal A$ is sync-maximal.
	By \cref{prelim:thm:SCC_sync_words}, $S$ is a sink component.
	By \crefitem{sync_max:cor:char_sync_max_aut}{cha_sync_max:item:Q3},
	every $2$-set is reachable in $\mathcal A$.
	So, we have $|Q \setminus S| \le 1$ by \cref{syncmax:prop:uppersets-num-sinks-size-Q}.
	If $Q = S$, then~$\mathcal A$ is strongly connected.
	So, by \cref{sync_max:cor:SCC}, it is completely reachable
	and completely distinguishable.
	Otherwise, let $q \in Q \setminus S$ and set $\Gamma = \{ x \in \Sigma \mid \delta(q, x) = q \}$.
	We have $Q = S \cup \{q\}$.
	As $S$ is reachable in~$\mathcal A$, by \cref{sync_maximal:lem:n-1-upperset} there
	exists $a \in \Sigma$ such that $\delta(Q, a) = S$
	and as $q \notin S$ we have $a \notin \Gamma$.
	Let $A \subseteq S$ be a non-empty
	subset. By \crefitem{sync_max:cor:char_sync_max_aut}{cha_sync_max:item:Q3}
	there exists $u \in \Sigma^*$ such that
	$\applytrans{\delta}{Q}{u} = A \cup \{q\}$.
	We have $\applytrans{\delta}{S}{u} \subseteq S$.
	Thus $\applytrans{\delta}{q}{u} \ne q$
	would give $\applytrans{\delta}{q}{u} \in S$
	and then $\delta(Q, u) \subseteq S$
	so that $q \notin \delta(Q, u)$.
	Hence, since $q \in \delta(Q, u)$, we must have $\applytrans{\delta}{q}{u} = q$ and 
	thus $u \in \Gamma^*$.
	This gives $\applytrans{\delta}{S}{u} = A$, and so 
	$\mathcal A_{|S, \Gamma}$ is completely reachable. 
	Complete distinguishability transfers from~$\mathcal A$
	to~$\mathcal A_{|S}$.
	Furthermore, all pairs of distinct subsets sharing the common state $q$
	are distinguishable, in particular those mentioned in Item~\ref{sync_maximal:thm:structure:item:dupl-set}.
		
	Conversely, if $\mathcal A$ is completely reachable and completely
	distinguishable, then it is sync-maximal by \cref{sync_max:cor:SCC}.
	Next, suppose $Q = S \cup \{q\}$ with $q \notin S$ and that,
	for $\Gamma = \{ x \in \Sigma \mid \delta(q, x) = q \}$, we have
	\begin{enumerate}
	\item $\delta(Q, a) = S$ for some $a \in \Sigma \setminus \Gamma$,
	
	\item for each $p \in Q$, all pairs of non-empty, non-singleton subsets
	of $\delta^{-1}(\dupl{a}, a) \cup \{p\}$ that are both distinct from
	$\delta^{-1}(\dupl{a}, a)$ and contain~$q$ are distinguishable,
	
	\item $\mathcal A_{|S}$ is completely
	distinguishable and~$\restrictAutGamS{\mathcal A}{\Gamma}{S}$ is completely reachable.
	\end{enumerate}
	Let $A \subseteq Q$ with $|A| \ge 2$.
	As $\restrictAutGamS{\mathcal A}{\Gamma}{S}$ is completely reachable,
	there exists $u \in \Gamma^*$ such that $\applytrans{\delta}{S}{u} = A \setminus \{q\}$.
	Then we have
	\[
	A = \begin{cases}
		\delta(Q, au) & \text{if } q \notin A, \\
		\delta(Q, u)  & \text{if } q \in A.		
	\end{cases}
	\]	
	So, every subset of size at least two is reachable in $\mathcal A$.    
    
    Now, let $A, B \subseteq Q$ be distinct with $|A|,|B| \ge 2$.    
    We can suppose $|A \cap S| \le 1$ or $|B \cap S| \le 1$,
    for if this is not the case,
    then, since $\mathcal A_{|S}$ is completely distinguishable,
    there exists a word $u$ mapping precisely one 
    of $A\cap S$ or $B \cap S$ to a singleton set.   
    In this case, if 
    $u$ does not distinguish $A$ and $B$, 
    we may argue instead with $\delta(A, u)$ and~$\delta(B, u)$ in what follows.
    So, without loss of generality, suppose $A = \{p, q\}$ for some $p \in S$.
    Then we have the following two cases.
    
    \begin{description}
    \item[$q \notin B$:] Then $B \subseteq S$. 
    Since $\restrictAutGamS{\mathcal A}{\Gamma}{S}$ is synchronizing,
    there exists $u \in \Sigma^*$ such that $|\delta(B, u)| = 1$.
        Then  $\delta(A, u) = \{\delta(p, u), q\}$
        with $\delta(p, u) \in S$ since $p \in S$.
        We have $q \notin S$ and so $u$ distinguishes $A$ and~$B$.
        
    \item[$q \in B$:] Suppose $\delta(A, a) = \delta(B, a)$.
    Since $A \ne B$ and $|B| \ge 2$,
    we have $\delta(s, a) \in \{\delta(p, a), \delta(q, a)\}$
    for some $s \in B \setminus A$. 
    The letter~$a$ has defect one and so there exists
    a unique pair of states that is mapped to the same state.
    This pair must be $s$ and $p$, or $s$ and~$q$,
    and with $A = \{p,q\}$ we conclude that $B \subseteq \{s, p, q\}$.
    Since $q \in B$ and $B \ne \{p,q\}$,
    we have $B = \{s,p,q\}$ or $B = \{s, q\}$.
    In the latter case, since $a$ does not map~$B$ to a singleton,
    we have $\delta(s, a) \ne \delta(q, a)$.    
    In all cases, neither $A$ nor $B$ equal $\delta^{-1}(\dupl{a}, a)$
    and both sets contain~$q$. By assumption, $A$ and~$B$ are distinguishable.    
        
    Next, suppose $\delta(A, a) \ne \delta(B, a)$.
    Since~$a$ has defect one, it cannot map both
    subsets to a singleton. Since $\mathcal A_{|S}$ is completely distinguishable,
    there exists a word $v$ distinguishing $\delta(A, a)$
    and~$\delta(B, a)$. Then $av$ distinguishes $A$ and~$B$.
    \end{description}
    
    Thus, $\mathcal A$ is completely distinguishable.    
    By \crefitem{sync_max:cor:char_sync_max_aut}{cha_sync_max:item:Q3}
    $\mathcal A$ is sync-maximal.
\end{proof}

The condition \crefitem{sync_maximal:thm:structure}{sync_maximal:thm:structure:item:AB} 
is rather technical. Next, \cref{sync_maximal:cor:structure,syncmax:prop:sufficient-dist-condition} 
give conditions that imply it.

\begin{corollary} \label{sync_maximal:cor:structure}
	Let $\semiautomat{A}{Q}{\Sigma}{\delta}$ and $S \subseteq Q$
	be a proper upper set such that $Q = S \cup \{q\}$ for some $q \in Q$.
	Set $\Gamma = \{ x \in \Sigma \mid \delta(q, x) = q \}$.
	Suppose the following holds:
	\begin{enumerate}
	\item We have $\delta(Q,a) = S$ for some $a \in \Sigma$, and
	
    \item \label{sync_maximal:cor:structure:2} for distinct $A,B \subseteq Q$ with $2 \le |A| \le 3$, $|B| = 2$
    and $q \in A \cap B$
    there exists $u \in \Sigma^*$ such that  $\delta(A, u) \ne \delta(B, u)$,
    $|\delta(A, u)| \ge 2$ or $|\delta(B, u)| \ge 2$,
    and $\delta(A, u) \subseteq S$
    or $\delta(B, u) \subseteq S$.	
	\end{enumerate}
	Then $\mathcal A$ is sync-maximal if and only if 
	$\mathcal A_{|S}$ is completely distinguishable
	and $\restrictAutGamS{\mathcal A}{\Gamma}{S}$ is completely reachable.
\end{corollary}
\begin{proof}
	By \cref{sync_max:cor:char_sync_max_aut} we can assume $|Q| \ge 3$.
	The set~$S$ is the unique sink component in $\mathcal A$.
	Thus, it equals the set of maximal states.
	
	First, suppose $\mathcal A$ is sync-maximal.
	By \cref{prelim:thm:SCC_sync_words}, $S$ equals the set of synchronizing states.
	By \cref{sync_maximal:thm:structure}, $\mathcal A_{|S}$ is completely distinguishable
	and $\restrictAutGamS{\mathcal A}{\Gamma}{S}$ is completely reachable.
	
	Conversely, suppose $\mathcal A_{|S}$ is completely distinguishable
	and $\restrictAutGamS{\mathcal A}{\Gamma}{S}$ is completely reachable.
	Then~$\mathcal A$ is synchronizing and, by \cref{prelim:thm:SCC_sync_words}, $S$
	is the set of synchronizing states.	
	Let $p \in S$ and $A,B \subseteq \delta^{-1}(\dupl{a}, a) \cup \{p\}$
	be distinct with $|A|, |B| \ge 2$ and $q \in A \cap B$.
	By assumption, we have $\delta(A, u) \ne \delta(B, u)$
	such that, without loss of generality, $\delta(A, u) \subseteq S$
	for some $u \in \Sigma^*$.
	If $|\delta(A, u)| = 1$ or $|\delta(B, u)| = 1$, 
	then $u$ distinguishes $A$ and~$B$.
	Otherwise, for $\delta(B, u) \subseteq S$ there exists a distinguishing word $v \in \Sigma^*$
	for $\delta(A, u)$ and $\delta(B, u)$. Then $uv$ distinguishes $A$ and~$B$.
	For $q \in \delta(B, u)$, we find $w \in \Gamma^*$
	such that $|\delta(A, uw)| = 1$ since $\restrictAutGamS{\mathcal A}{\Gamma}{S}$ is completely reachable.
	We have $\delta(q, w) = q$
	and $\delta(\delta(B, u) \cap S, w) \ne \emptyset$.
	This yields $|\delta(B, uw)| \ge 2$.
	Then, $uw$ distinguishes $A$ and~$B$.
	Thus, \crefitem{sync_maximal:thm:structure}{sync_maximal:thm:structure:item:AB} 
	holds and so we conclude that $\mathcal A$ is sync-maximal.
\end{proof}

Next, we list conditions that imply \cref{sync_maximal:cor:structure:2}
from \cref{sync_maximal:cor:structure} for specific sets.

\begin{proposition} \label{syncmax:prop:sufficient-dist-condition}
	Let $\semiautomat{A}{Q}{\Sigma}{\delta}$. Suppose $S \subseteq Q$
	is a proper upper set such that $\delta(Q,a) = S$ and 
	$Q = S \cup \{q\}$ for some $a \in \Sigma$ and some $q \in Q$.
    Set $\Gamma = \{ x \in \Sigma \mid \delta(q, x) = q \}$.
	
	\begin{enumerate}
	\item \label{syncmax:prop:sufficient-dist-condition:1} Let $A,B \subseteq Q$ be distinct with $|A|,|B| \ge 2$
	and $q \in A \cap B$.
	If $\delta(A, ua) \ne \delta(B, ua)$ for some $u \in \Gamma^*$,
	then $|\delta(A, ua)| \ge 2$ or $|\delta(B, ua)| \ge 2$.	
	
	\item \label{syncmax:prop:sufficient-dist-condition:2} If $\delta(S, a) = S$, then $\delta(A, a) \ne \delta(B, a)$
	for all distinct $A,B \subseteq Q$ with $|A|,|B| \ge 2$ 
	and $q \in A \cap B$.
	\end{enumerate}	
\end{proposition}
\begin{proof}
\ref{syncmax:prop:sufficient-dist-condition:1}: The set~$S$ is an upper set and 
 the sets $A \cap S$ and $B\cap S$ are non-empty. Thus, we have that $\delta(A, u) \cap S$
 and $\delta(B, u) \cap S$ are non-empty.
 We also have $q \in \delta(A, u) \cap \delta(B, u)$
 since $\delta(q, u) = q$.
 Thus $|\delta(A, u)| \ge 2$ and $|\delta(B, u)| \ge 2$.
 We have $\delta(A, u) \ne \delta(B, u)$ since $\delta(A, ua) \ne \delta(B, ua)$.
 Let $p \in \delta(A, u) \cap S$
 and $r \in \delta(B, u) \cap S$ with $p \ne r$.
 If $|\delta(A, ua)| = |\delta(B, ua)| = 1$,
 then $\delta(p, a) = \delta(q, a)$
 and $\delta(r, a) = \delta(q, a)$, 
 which is not possible as $a$ has defect one.
 Thus, $a$ can map at most one
 of both subsets $\delta(A, u)$
 or $\delta(B, u)$ to a singleton set, i.e.,
 we have $|\delta(A, ua)| \ge 2$ or $|\delta(B, ua)| \ge 2$.	

\ref{syncmax:prop:sufficient-dist-condition:2}: The letter~$a$ acts injectively on~$S$. Thus, we must have $q \in \delta^{-1}(\dupl{a}, a)$
and $\delta(A \cap S, a) \ne \delta(B \cap S, a)$
where $A\cap S$ and $B \cap S$ are non-empty.
In either case, whether $\delta(q, a) \in \delta(A\cap S, a) \cup \delta(B \cap S, a)$
or not, we find $\delta(A, a) \ne \delta(B, a)$.
\end{proof}

\begin{corollary}
	Let $\semiautomat{A}{Q}{\Sigma}{\delta}$
	be an SI-automaton.
	Then $\mathcal A$ is sync-maximal if and only if
	it is completely reachable and completely distinguishable
	or we have $Q = S \cup \{q\}$ for some upper set $S \subseteq Q$
	and $q \in Q \setminus S$ such that
	$\delta(Q, a) = S$ for some $a \in \Sigma$
	and~$\restrictAutGamS{\mathcal A}{\Gamma}{S}$ is completely reachable and completely
	distinguishable for $\Gamma = \{ x \in \Sigma \mid \delta(q, x) = q \}$.
\end{corollary}
\begin{proof}
	First, note that $\mathcal A_{|S}$
	is completely distinguishable
	if and only if $\restrictAutGamS{\mathcal A}{\Gamma}{S}$ is completely
	distinguishable since those letters in $\Sigma \setminus \Gamma$
	only move $q$ and fix all other states.	
	
	The letter~$a$ is a simple idempotent with $\delta(q, a) \ne q$.
	Thus $\delta(S, a) = S$.
	Then the claim follows 
	from \cref{sync_maximal:thm:structure,sync_maximal:cor:structure,syncmax:prop:sufficient-dist-condition}.
\end{proof}

\begin{corollary}
	Let $\semiautomat{A}{Q}{\Sigma}{\delta}$ with at least five states and $S \subseteq Q$
	be a proper upper set such that $\delta(Q, a) = S$ and $Q = S \cup \{q\}$ for 
	some $a \in \Sigma$ and some $q \in Q$.
	Suppose the set $S$ is cyclically permuted by some letter $b \in \Gamma$.
	Set $\Gamma = \{ x \in \Sigma \mid \delta(q, x) = q \}$.
	Then $\mathcal A$ is sync-maximal if and only if 
	$\mathcal A_{|S}$ is completely distinguishable
	and $\restrictAutGamS{\mathcal A}{\Gamma}{S}$ is completely reachable.
\end{corollary}
\begin{proof}
	Let $A, B \subseteq Q$ be distinct
	with $2 \le |A| \le 3$, $|B| = 2$,
	and $q \in A \cap B$.
	Set $T = \delta^{-1}(\dupl{a}, a)$.
	Suppose $\delta(A, a) = \delta(B, a)$.
	Then $A \setminus T = B \setminus T$
	and, since $A \ne B$, 
	we also have $A \cap T \ne \emptyset$ and $B \cap T \ne \emptyset$
	where $A \cap T \ne B \cap T$.		
	Set $n = |Q|$.
	Suppose $q \in T$. Let $T = \{p, q\}$.
	Then $A = \{s, p, q\}$ for some $s \in Q \setminus T$
	and $B = \{s, q\}$.
	Let $\delta(s, b^i) = p$ for some $0 < i < n$.
	If $i > 1$, set $u = b$
	and, otherwise, set $u = bb$.
	We have $\delta(A, u) \cap T = \{q\}$
	and $\delta(A, u) \setminus T \ne \delta(B, u) \setminus T$.
	Thus $\delta(A, ua) \ne \delta(B, ua)$.	
	Next, suppose $q \notin T$.
	Then $B = \{p, q\}$ for some $p \in T$.
	Let $u = b$ if $\delta(B, b) \cap T = \emptyset$
	or $u = bb$ otherwise. 
	Then $\delta(B, u) \cap T = \emptyset$.
	We have $\delta(A, u) \setminus T \ne \delta(B, u) \setminus T$.
	Thus $\delta(A, ua) \ne \delta(B, ua)$.	
	Since the image of~$a$ is~$S$,
	by \cref{sync_maximal:cor:structure,syncmax:prop:sufficient-dist-condition}
	we obtain the claim.
\end{proof}

Next, we list three corollaries to \cref{sync_maximal:thm:structure}. The first three follow readily,
the last one uses \cref{sync_max:cor:single-non-permutation}.

\begin{corollary}
	In a sync-maximal automaton $\mathcal A$
	with set of synchronizing states $S$,
	we have $|Q \setminus S| \le 1$
	and $\mathcal A_{|S}$
	is completely reachable
	and completely distinguishable.
\end{corollary}

\begin{corollary} \label{syncmax:cor:three-states-trap-not-syncmax}
A sync-maximal automaton with at least
three states has no trap state.
\end{corollary} 

\begin{corollary}
	A sync-maximal automaton is a $\rho$-automaton.
\end{corollary}

\begin{corollary}
	\label{sync_max:cor:structure_letters}
	Let $\semiautomat{A}{Q}{\Sigma}{\delta}$ be sync-maximal with $|Q| \ge 3$.
	If $\mathcal A$ is strongly connected, then $|\Sigma| \ge 2$
	where $\Sigma$ contains a letter of defect one
	and for $|\Sigma| = 2$ the other letter cyclically permutes the states. 
	If $\mathcal A$ is not strongly connected,
	then there exists a non-synchronizing state $q \in Q$
	such that $Q = S \cup \{q\}$ for the set of synchronizing states~$S$
	and with $\Gamma = \{ x \in \Sigma \mid \delta(q, x) = q \}$ we have:
	\begin{enumerate}
		\item \label{sync_max:cor:structure_letters:1} 
		each of the sets $\Sigma \setminus \Gamma$ and $\Gamma$ contains a letter of defect one,
				
		\item $|\Gamma| \ge 2$, $|\Sigma| \ge 3$,
		
		\item \label{sync_max:cor:structure_letters:3} if $|\Gamma| = 2$ and $|Q| \ge 4$, then
		$\Gamma$ also contains a permutational letter
		that cyclically permutes the states in~$S$ (this might not hold for $|Q| \le 3$, see \cref{sync_max:fig:ternary}). 
	\end{enumerate}
\end{corollary}
\begin{proof}
	First, assume $\mathcal A$ is strongly connected.
	By \cref{sync_max:cor:SCC}, $\mathcal A$
	is completely reachable.
	So, \cref{sync_max:cor2:single-non-permutation},
	we find that $|\Sigma| \ge 2$, $\Sigma$ contains a letter
	of defect one, and for $|\Sigma| = 2$ another letter cyclically permuting the states.
	Now, assume~$\mathcal A$ is not strongly connected.
	By \crefitem{sync_maximal:thm:structure}{sync_maximal:thm:structure:item:2nd}, 
	we have $Q = S \cup \{q\}$ for some state $q \in Q \setminus S$ where~$S$
	is the set of synchronizing states,
	that  $\Sigma \setminus \Gamma$ contains a letter of defect one,
	and that for $\Gamma = \{x\in \Sigma \mid \delta(q, x) = q \}$
	the automaton $\restrictAutGamS{\mathcal A}{\Gamma}{S}$
	is completely reachable.
	Applying \cref{sync_max:cor2:single-non-permutation}
	to the completely reachable automaton $\restrictAutGamS{\mathcal A}{\Gamma}{S}$
	yields the existence of a letter %$b \in \Gamma$ 
	in~$\Gamma$
	of defect one, $|\Gamma| \ge 2$, and so $|\Sigma| \ge 3$,
	and that for $|\Gamma| = 2$ and $|S| \ge 3$ the other letter in $\Gamma$ cyclically permutes~$S$.
\end{proof}

\section{Examples} 
\label{sync_max:sec:examples}

%!TEX root = ../../thesis-trier.tex


\Cref{sync_max:ex:compl-dist-non-cr,sync_max:ex:compl-dist-non-cr2,syncmax:ex:compl-dist-not-SCC} 
gave completely distinguishable automata that are not sync-maximal (which is easily seen with \cref{sync_maximal:thm:structure} 
in each case).
The automaton from \cref{sync_max:ex:2-sets-not-all-dist} 
is completely reachable, but not sync-maximal.
Here, \cref{ex:compl_reachable_not_sync_max} 
gives another completely reachable automaton that is not sync-maximal.
Then, \cref{ex:sync_max_not_compl_reachable_Cerny} gives two sync-maximal automata, derived from the \v{C}ern\'y automata,
that are not completely reachable.
\Cref{ex:three-state-ternary-sync-max} gives a three-state sync-maximal ternary automaton
without permutational letters, which is not possible
for more states by \cref{sync_max:cor:structure_letters}.
\Cref{sync_maximal:ex:structure:item:AB} gives an automaton that is not
sync-maximal but where all conditions from \crefitem{sync_maximal:thm:structure}{sync_maximal:thm:structure:item:2nd} except \crefitem{sync_maximal:thm:structure}{sync_maximal:thm:structure:item:AB} 
are fulfilled. 




\begin{example}
	\label{ex:compl_reachable_not_sync_max}
	The automaton (see Figure~\ref{fig:not_sync_max})
	with the state set $Q= \{1,2,3\}$, input letters $a, b, c, d$
	and transition function
	given by
	\begin{eqnarray*}
		\applytrans{\delta}{i}{a} = \begin{cases}
			2 &\text{if } i=1,2,\\
			3 &\text{if } i=3;
		\end{cases}
		&\qquad&
		\applytrans{\delta}{i}{b} =\begin{cases}
			1 &\text{if } i=1,2,\\
			3 &\text{if } i=3;
		\end{cases}
		\\
		\applytrans{\delta}{i}{c}=\begin{cases}
			1 &\text{if } i=1,2,\\
			2 &\text{if } i=3;
		\end{cases}
		&\qquad&
		\applytrans{\delta}{i}{d}=3\ \text{ for all }\ i=1,2,3.
	\end{eqnarray*}
	taken from~\cite[Ex.~2]{BondarV16} is completely reachable.
	However, it is not sync-maximal
	as, for example, the two $2$-sets $\{1,3\}$ 
	and $\{2,3\}$ are not distinguishable.
	\begin{figure}[ht]
		\begin{subfigure}[t]{0.31\textwidth}
			\centering 
		\scalebox{0.9}{%
			\begin{tikzpicture}[shorten >=1pt,->,node distance=2cm and 0.8cm]
				\node[state] (1) {1};
				\node[state, above right = of 1] (3) {3};
				\node[state, below right = of 3] (2) {2};
				
				\path[->]
				(1) edge [left,bend left] node {$d$} (3)
				(3) edge [bend left, right] node {$c$} (2)
				(2) edge [above] node {$b, c$} (1)
				(2) edge [bend left=5, left] node {$d$} (3)
				(1) edge [bend right,below] node {$a$} (2);
				
				\path[->]
				(1) edge [loop below] node {$b, c$} (1)
				(3) edge [loop above] node {$b, a, d$} (3)
				(2) edge [loop below] node {$a$} (2);
		\end{tikzpicture}}
			\caption{Non-sync-maximal completely reachable automaton from~\cite[Ex.~2]{BondarV16}.}
		\end{subfigure}\hfill%
		\begin{subfigure}[t]{0.31\textwidth}
			\centering 
			\scalebox{.9}{\begin{tikzpicture}[shorten >=1pt,->,node distance=2.5cm and 0.65cm]
					%\tikzstyle{state}=[circle,color=red!60,fill=black!10,minimum size=15pt,inner sep=1pt,very thick]
					%\tikzstyle{state}=[circle,draw,minimum size=15pt,inner sep=1pt,blur shadow,fill={rgb:black,0.4;white,10}]
					
					%\node (23) {$\{2,3\}$};
					%\node[below = of 23] (13) {$\{1,3\}$};
					%\node[below = of 13] (12) {$\{1,2\}$};			
					
					\node                  (123) {$\{1,2,3\}$};
					\node [below = of 123]  (13) {$\{1,3\}$};
					\node [right =  of 13]  (12) {$\{1,2\}$};
					\node [left = of 13]    (23) {$\{2,3\}$};
					\node [above of = 12]   (s) {$\{\ast \}$};
					
					\path[->] (123) edge [left]  node {$a$} (23)
					(123) edge [right] node {$b$} (13)
					(123) edge [right] node {$c$} (12);	
					
					\path[->] (12) edge node [align=left, right] {$a$ \\ $b$ \\ $c$} (s);			
					
					\path[->] (13) edge [loop below] node {$b$} (13);		
					
					\path[->] (23) edge [out=280,in=260,below] node {$c$} (12);
					\path[->] (23) edge [bend left,above]  node {$b$} (13);
					\path[->] (13) edge [bend left,below]  node {$a$} (23);
					\path[->] (13) edge [above] node {$c$} (12);
					\path[->] (23) edge [loop above] node {$a$} (23);		
			\end{tikzpicture}}%
			\caption{Power automaton, where singleton sets are merged, 
				the letter $d$, as it always ends in~$\{3\}$, and the empty set are omitted.} 
			%\caption{Part of the power automaton. Transitions into singleton subsets are omitted.}
			\label{fig:power-aut}
		\end{subfigure}\hfill%
		\begin{subfigure}[t]{0.31\textwidth}
			\centering 
			\scalebox{.9}{\begin{tikzpicture}[shorten >=1pt,->,node distance=1.75cm]
					%\tikzstyle{state}=[circle,color=red!60,fill=black!10,minimum size=15pt,inner sep=1pt,very thick]
					\tikzstyle{state}=[circle,draw,minimum size=15pt,inner sep=1pt,blur shadow,fill={rgb:black,0.4;white,10}]
					
					\node[state, initial] (1) {};
					\node[state, above of=1] (2) {};
					\node[state, above of=2, accepting] (3) {};
					
					\path[->] (1) edge [left] node {$c$} (2)
					(2) edge [left] node {$a, b, c, d$} (3)
					(1) edge [right, bend right] node {$d$} (3)
					(3) edge [loop above] node {$a, b, c, d$} (3)
					(1) edge [loop below] node {$a, b$} (1);
			\end{tikzpicture}}%
			\caption{Minimal automaton accepting set of synchronizing words, obtained
				from \cref{fig:power-aut} by merging $\{2,3\}$,$\{1,3\}$, and $\{1,2,3\}$.}
		\end{subfigure}	
		\caption[Completely reachable but not sync-maximal automaton]{
			Completely reachable automaton that is not sync-maximal~\cite[Ex.~2]{BondarV16}.}
		\label{fig:not_sync_max}
		%\caption{The automaton $\mathrsfs{E}_3$ and the graph $\Gamma_1(\mathrsfs{E}_3)$}\label{fig:e3}
	\end{figure}
	%The graph $\Gamma_1(\mathrsfs{E}_3)$ is shown in Fig.\,\ref{fig:e3} on the right; it is not \scn. However, it can be checked by a straightforward computation that the automaton $\mathrsfs{E}_3$ is completely reachable.
\end{example} 

\begin{example} \label{ex:sync_max_not_compl_reachable_Cerny}
	Let $\semiautomat{A}{Q}{\{a,b,c\}}{\delta}$
	be one of the automata from \cref{sync_max:fig:c-not-permute,sync_max:fig:A-n}
	with $n \ge 5$.
	The set $S = \{0,1,\ldots,n-1\}$ is an upper set
	and~$\restrictAutGamS{\mathcal A}{\{a,b\}}{S}$
	equals the \Cerny automaton,
	which is completely reachable and completely distinguishable	
	by \cref{sync_max:thm:cerny_sync_maximal}.
	Next, consider distinct $A,B \subseteq Q$ with $n \in A \cap B$.
	Without loss of generality, $q \in A \setminus B$ for some $q \in Q$.
	Then there exists $i \ge 0$ such that 
	$1 \in \delta(A, b^i) \setminus \delta(B, b^i)$
	and so $\delta(A, b^ic) \ne \delta(B, b^ic)$.
	Also, we have $\delta(Q, c) = S$.   		
	So, by \cref{sync_maximal:cor:structure} 
	we obtain that~$\mathcal A$
	is sync-maximal.	
	
	\begin{figure}[ht]
		\centering
		\begin{subfigure}[b]{0.34\textwidth}
			\centering
			\scalebox{0.9}{
				\begin{tikzpicture}[shorten >=1pt,->,node distance=5cm]
					%\tikzstyle{vertex}=[circle,color=red!60,fill=black!10,minimum size=25pt,inner sep=1pt,very thick]
					%\tikzstyle{vertex}=[circle,draw,minimum size=5pt,inner sep=1pt]
					
					\foreach \name/\angle/\text in {P-n-2/234/n-2, P-n-1/162/n-1, 
						P-0/90/0, P-1/18/1, P-2/-54/2}
					\node[state,blur shadow, fill={rgb:black,0.4;white,10}, inner sep=0pt, minimum size=30pt] (\name) at (\angle:1.75cm) {$\text$};  %{\small $\text$};
					
					\node[state,blur shadow,fill={rgb:black,0.4;white,10},left = 0.6cm of P-0, inner sep=0pt, minimum size=30pt,yshift=0.25cm] (n) {$n$};
					
					
					\path[->] %(n) edge [loop above] node {$a,b$} (n)
					(n) edge [above] node {$c$} (P-0);
					%	(P-0) edge [loop above] node {$c$} (P-0);
					
					\path[->] (P-n-2) edge [left] node {$b,c$} (P-n-1);
					\path[->] (P-n-1) edge [above,pos=0.35] node {$b$} (P-0);
					%\path[->] (P-n-1) edge [loop left] node {$a, c$} (P-n-1);
					%\path[->] (P-n-1) edge [loop left] node {$a,c$} (P-n-1);
					\path[->] (P-0) edge [right, pos=0.2,inner sep=11pt] node {$a, b, c$} (P-1);
					\path[->] (P-1) edge [right] node {$b,c$} (P-2);
					
					%\path[->] (P-n-2) edge [loop left] node {$a,c$} (P-n-2);
					%\path[->] (P-1) edge [loop right] node {$a,c$} (P-1);
					%	\path[->] (P-2) edge [loop right] node {$a,c$} (P-2);
					
					%\path[->] (n) edge [loop left] node {$a,b$} (n);
					\node at (0.1, -1.4) {\large $\ldots$};
					
					%	\path[->] (P-1)   edge [loop right]  node {$a$} (P-1)
					%	(P-n-2)   edge [loop left]  node {$a$} (P-n-2)
					%	(P-2)   edge [loop right] node {$a$} (P-2);
				\end{tikzpicture}
			}
			\caption{Sync-maximal, but not completely reachable.}
			\label{sync_max:fig:c-not-permute}
		\end{subfigure}
		\hfill
		\begin{subfigure}[b]{0.34\textwidth}
			\centering
			\scalebox{0.9}{
				\begin{tikzpicture}[shorten >=1pt,->,node distance=5cm]
					%\tikzstyle{vertex}=[circle,color=red!60,fill=black!10,minimum size=25pt,inner sep=1pt,very thick]
					%\tikzstyle{vertex}=[circle,draw,minimum size=5pt,inner sep=1pt]
					
					\foreach \name/\angle/\text in {P-n-2/234/n-2, P-n-1/162/n-1, 
						P-0/90/0, P-1/18/1, P-2/-54/2}
					\node[state,blur shadow, fill={rgb:black,0.4;white,10}, inner sep=0pt, minimum size=30pt] (\name) at (\angle:1.75cm) {$\text$};  %{\small $\text$};
					
					\node[state,blur shadow,fill={rgb:black,0.4;white,10},left = 0.6cm of P-0, inner sep=0pt, minimum size=30pt,yshift=0.25cm] (n) {$n$};
					
					\path[->] %(n) edge [loop above] node {$a,b$} (n)
					(n) edge [above] node {$c$} (P-0);
					%	(P-0) edge [loop above] node {$c$} (P-0);
					
					\path[->] (P-n-2) edge [left] node {$b$} (P-n-1);
					\path[->] (P-n-1) edge [above,pos=0.35] node {$b$} (P-0);
					%\path[->] (P-n-1) edge [loop left] node {$a,c$} (P-n-1);
					\path[->] (P-0) edge [above] node {$a, b$} (P-1);
					\path[->] (P-1) edge [right] node {$b$} (P-2);
					
					%\path[->] (P-n-2) edge [loop left] node {$a,c$} (P-n-2);
					%\path[->] (P-1) edge [loop right] node {$a,c$} (P-1);
					%	\path[->] (P-2) edge [loop right] node {$a,c$} (P-2);
					
					\node at (0.1, -1.4) {\large $\ldots$};
				\end{tikzpicture}
			}
			\caption{Sync-maximal, but not completely reachable.}
			\label{sync_max:fig:A-n}
		\end{subfigure}
		\hfill
		\begin{subfigure}[b]{0.25\textwidth}
			\centering
			\scalebox{0.9}{
				\begin{tikzpicture}[shorten >=1pt,->] %,node distance=2cm]
					\node[state] (0) {$0$};
					\node[state, below = of 0] (1) {$1$};
					\node[state, right = of 1] (2) {$2$};
					
					\path[->] (0) edge [left]  node {$a$} (1) 
					(1) edge [bend left, above] node {$b$} (2)
					(2) edge [bend left, below] node {$c$} (1);
				\end{tikzpicture}
			}
			\caption{Sync-maximal automaton with no permutational letters.}
			\label{sync_max:fig:ternary}
		\end{subfigure}
		\hfill
		%		\begin{subfigure}[b]{0.45\textwidth}
			%			\centering
			%			\scalebox{0.9}{
				%				\begin{tikzpicture}
					%					\node[state] (1) {$1$};
					%					\node[state, right = of 1] (2) {$2$};
					%					\node[state, right = of 2] (3) {$3$};
					%					
					%					\path[->] (1) edge [bend left,above] node {$a$} (2)
					%					(2) edge [bend left,above] node {$b$} (3)
					%					(3) edge [bend left,below] node {$c$} (2)
					%					(2) edge [bend left,below] node {$d$} (1);
					%					%	(1) edge [loop above] node {$e$} (1)
					%					%	(2) edge [loop above] node {$e$} (2)
					%					%	(3) edge [loop above] node {$e$} (3);
					%				\end{tikzpicture}
				%			}
			%			\caption[Pairs of $2$-sets distinguishable but not completely distinguishable]{Completely reachable, not completely distinguishable.}	
			%			\label{sync_max:fig:SCC_not_sufficient}
			%		\end{subfigure}
		\caption[Sync-maximal automata but not completely reachable automata]{Different sync-maximal automata that are not completely reachable, self-loops omitted. \Cref{sync_max:fig:c-not-permute,sync_max:fig:A-n,sync_max:fig:A-n}
			are derived from the \v{C}ern\'y automata by adding a ``dangling'' state,
			see \cref{ex:sync_max_not_compl_reachable_Cerny}.}		
		\label{fig:three graphs}
	\end{figure}	
\end{example}

\begin{example}\label{ex:three-state-ternary-sync-max}
	By \crefitem{sync_max:cor:structure_letters}{sync_max:cor:structure_letters:3}
	a ternary sync-maximal automaton with at least four states has a permutational letter.
	\Cref{sync_max:fig:ternary} gives a three-state sync-maximal automaton over a ternary alphabet
	with no permutational letter.
\end{example}

\begin{example} \label{sync_maximal:ex:structure:item:AB}
	Here, we construct automata that are not strongly connected and where all conditions
	from \crefitem{sync_maximal:thm:structure}{sync_maximal:thm:structure:item:AB}
	except \crefitem{sync_maximal:thm:structure}{sync_maximal:thm:structure:item:AB}
	are fulfilled; \cref{sync_maximal:fig:structure:item:AB} illustrates the construction in the case of four states.	
	Let $\semiautomat{A}{Q}{\Sigma}{\delta}$ with $Q = \{0,1,\ldots,n\}$ for $n \ge 3$ and where, for $S = Q \setminus \{0\}$,	we have $\Sigma = \{a\} \cup \Gamma$ with $\Gamma = \{ b_A \mid \emptyset \ne A \subsetneq S \} \cup \{ c_q \mid q \in S \setminus \{1\}\}$ and
	\begin{enumerate}
		\item $\delta(q, a) = \min(q+1,n)$ for $q \in Q$,
		\item for $A \subseteq S$ with $A = \{p_1, \ldots, p_m\}$ where $p_1 < \ldots < p_m$ and $1 \le m < n$
		we have $\delta(q, b_A) = p_{\min(q, m)}$ for $q \in S$
		and $\delta(0, b_A) = 0$,
		\item for $q \in S \setminus \{1\}$ and $p \in Q$ we have
		\[
		\delta(p, c_q)
		=
		\begin{cases}
			0     & \text{if } p = 0, \\
			n     & \text{if } p \notin \{0,q\}, \\
			n - 1 & \text{if } p = q. 	
		\end{cases}
		\]
	\end{enumerate}
	%Figure~\ref{sync_maximal:fig:structure:item:AB} depicts $\mathcal A$ for $Q = \{0,1,2,3\}$. 
	Now, the letters in $\Gamma$ are precisely those fixing~$0$, we have	$\delta(Q, a) = S$, $S$ equals precisely the set of synchronizing states and $\mathcal A_{|\Gamma, S}$ is completely reachable as for each $A \subseteq Q$ with $1 \le |A| < n$ we have $\delta(S, b_A) = A$. Furthermore, as $\delta(S \setminus \{1\}, b_{\{1,2\}}) = \{2\}$ and $\delta(S, b_{\{1,2\}}) = \{1,2\}$, and $\delta(S \setminus \{q\}, c_q) = \{n\}$	and $\delta(S, c_q) = \{n, n-1\}$ for $q \in \{2,\ldots,n\}$, 
	$\restrictAutGamS{\mathcal A}{\Gamma}{S}$ is completely distinguishable	by Theorem~\ref{sync_max:thm:char_distinguishable}~\ref{dist:item:single_point}. Next, we have	$\delta(\{0,n-1\}, c_n) = \{0,n\}$ and $\delta(\{0,n\}, c_n)  = \{0,n-1\}$, $\delta(\{0,n-1\}, c_{n-1})  = \{0,n-1\}$ and $\delta(\{0,n\}, c_{n-1})  = \{0,n\}$, for $q \notin \{n-1,n\}$
	we have $\delta(\{0,n-1\}, c_q) = \delta(\{0,n\}, c_q) = \{0,n\}$, $\delta(\{0,n-1\}, a)  = \delta(\{0,n\}, a) = \{1,n\}$	and, for $\emptyset \ne A \subsetneq S$, we have
	\[
	\delta(\{0,n-1\}, b_A) = \delta(\{0,n\}, b_A) = \{0,\max(A)\}.
	\]
	So, $\{0,n-1\}$ and $\{0,n\}$ are not distinguishable in~$\mathcal A$ and \crefitem{sync_maximal:thm:structure}{sync_maximal:thm:structure:item:AB} is the only condition from \cref{sync_maximal:thm:structure} that fails.	
	
	\begin{figure}\centering 
		\scalebox{0.95}{
			\begin{tikzpicture}[shorten >=1pt,->,node distance=2cm and 5cm,>={Stealth[length=2mm]}]
				\node[state] (0) {$0$};
				\node[state, yshift=1cm, below = of 0] (1) {$1$};	
				\node[state, above right = of 1] (2) {$2$};
				\node[state, below right = of 2] (3) {$3$};
				
				\path[->] (0) edge [left] node {$a$} (1)
				(1) edge [bend left=5, above, sloped] node {$a, b_{\{2\}}, b_{\{2,3\}}$} (2)
				(2) edge [bend left=5, above, sloped] node {$a, b_{\{3\}}, b_{\{1,3\}}, b_{\{2,3\}}, c_3$} (3);
				
				\path[->] (3) edge [bend left=5, below, sloped] node {$b_{\{2\}}, b_{\{1,2\}}, c_3$} (2)
				(2) edge [bend left=5, below, sloped] node {$b_{\{1\}}$} (1);
				
				\path[->] (1) edge [bend left=5, above] node {$b_{\{3\}}, c_2, c_3$} (3)
				(3) edge [bend left=5, below] node {$b_{\{1\}}$} (1);
		\end{tikzpicture}}
		\caption[One property from \cref{sync_maximal:thm:structure} fails and so not sync-maximal]{Automaton  (self-loops omitted) according to \cref{sync_maximal:ex:structure:item:AB}.}
		\label{sync_maximal:fig:structure:item:AB}
	\end{figure}
\end{example} 

\section{Discussion}


 \Cref{sync_maximal:thm:structure} and \cref{sync_max:cor:structure_letters}
 gave structural results linking the class of sync-maximal automata
 to completely reachable and completely distinguishable automata.
 Later, in \cref{sync_max:sec:SI}, we will give more precise
 structural results for classes of automata.
 For example, an automaton whose transformation semigroup contains
 a transitive group of units and where every letter has defect at most one is synchronizing
 if and only if it is sync-maximal, see \cref{sync_max:sec:defectone}.
 This statement applies to the the \Cerny automata from \cref{intro:fig:cerny}.
 
We note that deciding sync-maximality is \PSPACE-complete, see
the discussion in \cref{sync_max:decision:discussion}.

\IN{reset complexity}
By \cref{prelim:thm:sync_words_is_ideal}, languages of the form $\Syn(\mathcal A)$ are precisely the regular ideal languages.
Thus, shortest words in regular ideal languages corresponds
to shortest paths in accepting automata for sets of synchronizing words.
Now, investigating the length of the shortest words in regular ideal languages as a language-theoretic approach to the \v{C}ern\'{y} conjecture has been pursued and extended by several authors~\cite{DBLP:journals/ijfcs/AlmeidaR16,DBLP:journals/jalc/FrigeriR25,DBLP:journals/jalc/FrigeriR19,DBLP:journals/fuin/MaslennikovaR18,DBLP:journals/iandc/PribavkinaR11,ReisR16,DBLP:journals/aam/Rodaro18}, based on a method introduced by Rodero et al.~\cite{DBLP:journals/fuin/MaslennikovaR18,ReisR16,DBLP:journals/aam/Rodaro18} that even associates a strongly connected synchronizing automaton with each regular ideal language.
Maslennikova~\cite{DBLP:journals/corr/Maslennikova14,DBLP:journals/ijfcs/Maslennikova19} 
introduced the \emph{reset complexity} 
of a regular ideal language $I \subseteq \Sigma^*$ 
as the number of states of a smallest automaton~$\mathcal A$ 
such that $I = \Syn(\mathcal A)$.
For a sync-maximal automaton $\mathcal A$,
the gap between the size of an accepting automaton for $\Syn(\mathcal A)$, i.e., the state complexity of this set,
and the size of an automaton having $\Syn(\mathcal A)$ as its set of synchronizing words, i.e., its reset complexity,
is maximal possible (note that in this case, no automaton with fewer states than~$\mathcal A$ can have $\Syn(\mathcal A)$
as its set of synchronizing words).
In this sense, for sync-maximal automata the complexity of the set of synchronizing words is maximal.